\section{Introduction}
Day-ahead prices describe the hourly cost of electricity for the next delivery
day. A forecast can be accurate on average yet miss by a large margin under
particular price or weather conditions. An annual error score cannot identify
those conditions or show whether they persist. We address two questions:

\noindent\textbf{RQ1.} Which combinations of calendar conditions and market
regimes are associated with unusually large errors in DK1 day-ahead forecasts?

\noindent\textbf{RQ2.} Do the associations persist in a later period and across
different forecast models?

We generate comparable errors from chronological forecasts, fit K-means to
summarize the market conditions of each day, and use Apriori to find interpretable
combinations of conditions. Temporal evaluation tests whether the rules persist
in the following year. Comparing the same rules across a sparse linear model,
tree ensembles, and a neural network can reveal whether they reflect a weakness
specific to the Ridge baseline.
